\section{Numbers, quantities, and units}

A number by itself is often not enough. In science, we usually work with quantities. A quantity has a numerical value, a unit, and a meaning.

\begin{center}
\begin{tabular}{@{}lll@{}}
\toprule
Expression & What it contains & Interpretation \\
\midrule
\(5\) & number only & incomplete information \\
\(5\,\mathrm{m}\) & value and unit & length \\
\(5\,\mathrm{s}\) & value and unit & time \\
\(5\,\mathrm{kg}\) & value and unit & mass \\
\(5\,\mathrm{C}\) & value and unit & electric charge \\
\bottomrule
\end{tabular}
\end{center}

These expressions do not mean the same thing. The unit is part of the information.

\begin{remarkbox}
Units are not decoration. They are part of the quantity and must be carried through a calculation consistently.
\end{remarkbox}

We can add quantities with the same unit:
\[
2\,\mathrm{m} + 3\,\mathrm{m} = 5\,\mathrm{m}.
\]
But the expression
\[
2\,\mathrm{m} + 3\,\mathrm{s}
\]
is not a meaningful sum, because metres and seconds describe different kinds of quantities.

Units also help us check formulas. If the left-hand side of an equation has a certain unit, then the right-hand side must have the same unit. This simple observation is one of the most useful ways to detect errors.

\begin{summarybox}
A calculation may be numerically correct and still be meaningless if the units do not match. Checking units does not solve the problem, but it protects us from many invalid results.
\end{summarybox}
